\section{Introdução}

Na AGES IV, no segundo semestre de 2026, passei a integrar o projeto ReadRace, sob orientação do professor Gabriel Fonseca Silva. O projeto propõe um aplicativo mobile para incentivar o hábito da leitura por meio de biblioteca pessoal, acompanhamento de progresso, interação social e gamificação. Seus stakeholders são Bernardo Schedler, Gabriela de Moraes, José Pedro Schury e Monique Soares Silveira.

O problema apresentado pelos clientes envolve a dificuldade de manter a constância na leitura diante das distrações digitais. A proposta é tornar esse hábito mais estimulante por meio de metas, desafios e recompensas. Conforme o escopo refinado pela equipe, a experiência deve valorizar o progresso sem punir os períodos de ausência. A leitura dos livros ocorre fora da plataforma: o aplicativo acompanha a atividade do leitor, sem hospedar o conteúdo das obras. Pagamentos e monetização também estão fora do escopo.

Nesta etapa, minha atuação passou a envolver mais diretamente a gestão do projeto, em conjunto com os demais colegas de AGES IV. Participei de alinhamentos, preparação de perguntas aos clientes, discussão de user stories, revisão da modelagem de dados, acompanhamento de tarefas e apresentações. Também contribuí com desenvolvimento e revisão de código, especialmente na aproximação da primeira entrega.

Este capítulo apresenta um relato parcial da AGES IV, abrangendo a preparação do projeto e a Sprint 1, cuja entrega aos stakeholders ocorreu em 18 de setembro de 2026. As ações definidas para a Sprint 2 são mencionadas somente como próximos passos. Os registros pessoais utilizados abrangem 13 de agosto a 17 de setembro de 2026. A organização distingue a preparação da execução, sem estabelecer datas formais de início e fim da Sprint 0 que não foram confirmadas.
